% TOP_PROBLEM: {"id":30007669,"problem_number":"LOCAL-30007669","title":"Information that creates electoral accountability","category_id":21,"category":{"id":21,"name":"economics","display_name":"Economics","description":"Open economic theory statements, published research agendas, and proposed scoped specifications; question provenance and historical priority are qualified per record.","slug":"economics","order_index":21,"created_at":"2026-10-08T00:00:00Z"},"status":"provenance_pending","external_url":null,"legacy_ids":[],"published":false,"statement":"In the explicitly specified setting: When does credible information about politician performance change voting and public-service delivery, and when do identity, distrust, clientelism, or voter coordination overwhelm it? The mathematical statement above fixes the target, assumptions and scope of this catalogue formulation.","economics":{"original_id":158,"field":"Political economy and institutions","kind":"identification","formalization_basis":"proposed_specification","source_status":"not_verified","priority_status":"not_established","priority_note":"No qualifying problem statement was directly verified, so historical priority cannot be assigned.","earliest_verified_reference":null,"current_reference":{"authors":["Thad Dunning","Guy Grossman","Macartan Humphreys","Susan D. Hyde","Craig McIntosh","Gareth Nellis","Claire L. Adida","Eric Arias","Clara Bicalho","Taylor C. Boas","Mark T. Buntaine","Simon Chauchard","Anirvan Chowdhury","Jessica Gottlieb","F. Daniel Hidalgo","Marcus Holmlund","Ryan Jablonski","Eric Kramon","Horacio Larreguy","Malte Lierl","John Marshall","Gwyneth McClendon","Marcus A. Melo","Daniel L. Nielson","Paula M. Pickering","Melina R. Platas","Pablo Querubín","Pia Raffler","Neelanjan Sircar"],"title":"Voter information campaigns and political accountability: Cumulative findings from a preregistered meta-analysis of coordinated trials","year":2019,"url":"https://pmc.ncbi.nlm.nih.gov/articles/PMC6609214/","doi":"10.1126/sciadv.aaw2612","locator":"Abstract, paragraphs 1–3; Introduction, paragraphs 4–6","evidence_summary":"The coordinated trials provide null average electoral effects and exploratory conditions; no explicit remaining long-run services question was verified."},"formalization_note":"This is a proposed scoped research specification. No primary passage explicitly stating this entire remaining question as open was verified. The supporting literature may answer important subquestions; this entry must not be advertised as a verified formal open problem."},"statusQualification":"Provenance pending: an explicit published statement of this exact open question was not verified. This is a catalogue-authored candidate specification, not a certified open conjecture. This is a proposed scoped research specification. No primary passage explicitly stating this entire remaining question as open was verified. The supporting literature may answer important subquestions; this entry must not be advertised as a verified formal open problem.","statusReviewedAt":"2026-10-08"}
% FIRST_POSTED: 2026-10-08
% ENTRY_KIND: statement_only
% !TeX program = lualatex
\documentclass[11pt]{article}
\usepackage{fontspec}
\usepackage[a4paper,margin=25mm]{geometry}
\usepackage{amsmath,amssymb,mathtools}
\usepackage{xurl}
\usepackage[hidelinks]{hyperref}
\newcommand{\sref}[2]{\hyperlink{source:#1:#2}{[S#2]}}
\setlength{\emergencystretch}{3em}
\begin{document}
% problemId:30007669
\section*{Information that creates electoral accountability}\label{30007669}
\subsection{Definitions and mathematical statement}
Fix an electoral setting and a verified measure \(q_j\) of incumbent \(j\)'s public-service performance. An information intervention \(z\) delivers a declared message through a specified channel before an election. Let \(B_i(z)\) be voter \(i\)'s belief about performance, \(V_i(z)\) the vote choice, and \(G_{jT}(z)\) subsequent service provision. Define belief, vote, and service effects separately, including \(\tau_G=E[G_{jT}(1)-G_{jT}(0)]\). Accountability additionally requires evaluating how support responds to performance, rather than an undirected change in incumbent vote share. Orthogonal message, source-credibility, coordination, or alternative-candidate variations can distinguish information failures from distrust, identity preferences, clientelistic incentives, and inability to coordinate. State assumptions permitting those mechanism interpretations. Account for message spillovers, selective exposure, turnout, and politicians' responses to anticipated disclosure. A null average vote effect can coexist with corrected beliefs or subgroup effects; it does not show that information never matters. The proposed problem is to identify conditions yielding durable accountability and service improvements. Existing coordinated trials provide substantial answers for their tested designs; an exact primary open-question paragraph for the combined catalogue formulation was not verified.

\par\textbf{Formulation and scope.} This is a proposed scoped research specification. No primary passage explicitly stating this entire remaining question as open was verified. The supporting literature may answer important subquestions; this entry must not be advertised as a verified formal open problem.

\par\textbf{Open-question provenance.} An explicit published open-question statement was not verified. Sources below are supporting literature and must not be treated as a first listing.

\par\textbf{Historical priority.} No qualifying problem statement was directly verified, so historical priority cannot be assigned.

\subsection{Short English statement}
In the explicitly specified setting: When does credible information about politician performance change voting and public-service delivery, and when do identity, distrust, clientelism, or voter coordination overwhelm it? The mathematical statement above fixes the target, assumptions and scope of this catalogue formulation.

\subsection{Sources}
\begin{itemize}\raggedright
\item[\textnormal{[S1]}]\hypertarget{source:30007669:1}{} Thad Dunning, Guy Grossman, Macartan Humphreys, Susan D. Hyde, Craig McIntosh, Gareth Nellis, Claire L. Adida, Eric Arias, Clara Bicalho, Taylor C. Boas, Mark T. Buntaine, Simon Chauchard, Anirvan Chowdhury, Jessica Gottlieb, F. Daniel Hidalgo, Marcus Holmlund, Ryan Jablonski, Eric Kramon, Horacio Larreguy, Malte Lierl, John Marshall, Gwyneth McClendon, Marcus A. Melo, Daniel L. Nielson, Paula M. Pickering, Melina R. Platas, Pablo Querubín, Pia Raffler, Neelanjan Sircar. 2019. Voter information campaigns and political accountability: Cumulative findings from a preregistered meta-analysis of coordinated trials. Abstract, paragraphs 1–3; Introduction, paragraphs 4–6.
\url{https://pmc.ncbi.nlm.nih.gov/articles/PMC6609214/}.
DOI: 10.1126/sciadv.aaw2612.
{\small\textit{Supporting or current literature; see evidence qualification. The coordinated trials provide null average electoral effects and exploratory conditions; no explicit remaining long-run services question was verified.}}
\end{itemize}
\end{document}
